\section{Introduction}
Experimental design is a classical and widely studied topic dating back to Fisher’s foundational work \citep{fisher1935design}, with applications spanning clinical trials, A/B testing, engineering, and scientific discovery. A major part of its appeal is that it allows one to commit in advance to a fixed set of measurements to collect, that is, a non-adaptive design; in many linear settings, there is a rich theory for constructing such designs \citep{pukelsheim2006optimal}.  Such non-adaptive designs also offer practical advantages over adaptive approaches, including computational simplicity and the ability to parallelize or batch data collection. Moreover, in certain linear settings, such non-adaptive approaches are known to be near-minimax optimal \citep{arias2012fundamental,even2002pac,fan2025universal}. Motivated by these advantages, we revisit a fundamental sequential decision problem in a linear setting, namely {pure exploration in} linear bandits, and ask how much adaptivity can help beyond the best non-adaptive design. In particular, we study the minimax complexity of non-adaptive designs and ask whether the structure of the action set confines adaptive sampling to logarithmic gains, or allows polynomial improvements over any non-adaptive design.



In this paper, we study the $\varepsilon$-best arm identification  problem {(i.e., pure exploration)} in linear bandits. We are given a compact arm set $\mc{X}\subset \R^d$ with $\mathrm{span}(\mc{X})=\R^d$ and an unknown reward vector $\theta\in\R^d$. At each round $t=1,2,\dots$, an algorithm selects an arm $x_t\in\mc{X}$ and observes $y_t=\langle x_t,\theta\rangle+\eta_t$, where $\eta_t\sim\mc{N}(0,1)$\footnote{All our algorithmic results extend to independent mean-zero subgaussian noise with bounded variance proxy. We restrict attention to Gaussian noise for simplicity of presentation.}.
 Given $\varepsilon\in(0,1)$ and $\delta\in(0,1)$, the objective is to design an algorithm, consisting of a sampling rule and a stopping time, such that upon stopping it outputs $\widehat{x}\in\mc{X}$ satisfying $\P\!\left(\max_{x\in\mc{X}}\langle x-\widehat{x},\theta\rangle\le \varepsilon\right)\ge 1-\delta$. We call an algorithm non-adaptive if its sampling rule does not depend on the observed rewards $y_t$, and adaptive otherwise. Our goal is to characterize the power and limitations of adaptivity in this problem setting.

When $\mathcal{X}=\{e_1,e_2,\ldots,e_d\}$, the problem reduces to the classical stochastic multi-armed bandit setting. A simple non-adaptive algorithm samples each arm $\mathcal{O}\!\left(\frac{\log(d/\delta)}{\varepsilon^2}\right)$ times and outputs the arm with the largest empirical mean. A union bound shows that, with probability at least $1-\delta$, the returned arm is $\varepsilon$-best. One way to extend this approach to linear bandits is via a fixed design. In optimal experimental design, a design distribution $\lambda$ is a probability distribution over $\mathcal{X}$. Using a finitely supported $\lambda$, one can construct a fixed design, that is, a deterministic sequence of arms, in various ways. Typically, the number of times an arm $x\in\mathcal{X}$ appears in the sequence is chosen proportional to $\lambda(x)$, though other constructions are also possible.
The sampling rule then pulls the arms in this predetermined order and observes the corresponding rewards. Using these observations, one typically forms an unbiased estimator $\widehat{\theta}$ and outputs the empirically best arm. We refer to such procedures as non-adaptive fixed-design algorithms.

One can view the non-adaptive multi-armed bandit algorithm discussed above as a fixed-design method induced by the uniform design distribution over the arms. This procedure is already near-optimal for multi-armed bandits: adaptive algorithms such as Median Elimination \citep{even2002pac} achieve the optimal minimax sample complexity $\mathcal{O}\!\left(\frac{d\log(1/\delta)}{\varepsilon^2}\right)$, improving over the optimal non-adaptive rate only by a logarithmic factor. This motivates the following question for the more general linear bandit setting:
\begin{quote}
\emph{What is the minimax sample complexity of non-adaptive fixed-design algorithms for $\varepsilon$-best arm identification in linear bandits? {Moreover, can adaptivity improve over the optimal non-adaptive rate only by logarithmic factors, as in multi-armed bandits, or are polynomial-factor improvements possible?}}
\end{quote}

\subsection{Notations}
Before we address the question raised above, we would like to mention a few notations. The unit $\ell_2$ ball refers to the set $\mathbb{B}_d := \{x \in \mathbb{R}^d : \|x\|_2 \le 1\}$.  If $\lambda$ is a distribution over the set $\mathcal{X}$, we say $\lambda \in \triangle_{\mathcal{X}}$ and will denote $A(\lambda; \mc{X}) := \mathbb{E}_{x\sim\lambda}\left[x x^\top\right]$. We define the gaussian width term $w(\mathcal{X})$ depending on $\mathcal{X}$ as 
\begin{align}
    w(\mathcal{X}):=\inf_{\lambda \in \triangle_{\mathcal{X}}}\mathbb{E}_{\eta\sim\mathcal{N}(0,I_d)}\left[\max_{x\in\mathcal{X}}\langle x,A(\lambda,\mc{X})^{-1/2}\eta \rangle\right]. \label{eqn:gauss_width}
\end{align}
For any matrix $W\in \mathbb{R}^{d\times d}$ and a vector $x\in\mathbb{R}^d$, we will denote $\|x\|_W^2 := x^\top W x$. 
For any positive scalar $x$ define $(x)_+ = \max\{1,x\}$ as the maximum of $x$ and one, \emph{not} zero.
Define $x_* = \arg\max_{x \in \mc{X}} \langle x ,\theta \rangle$.

In this paper, we focus on $(\varepsilon,\delta)$-PAC algorithms that, for any $\varepsilon\in(0,1]$ and $\delta\in(0,1)$, use $\frac{H_1\log(1/\delta)+H_2}{\varepsilon^2}$ samples, where $0<H_1\leq H_2$, and output $\hat{x}\in\mathcal{X}$ such that $\mathbb{P}(\langle x_*-\hat{x},\theta\rangle \le \varepsilon) \ge 1-\delta$.
\subsection{Our Contributions and Techniques}
We now address the question posed above by summarizing our main contributions and techniques. We first present a non-adaptive algorithm based on a fixed design with sample complexity $\mathcal{O}\!\left(\frac{d\log(1/\delta)}{\varepsilon^2}+\frac{w(\mathcal{X})^2}{\varepsilon^2}\right)$, where $w(\mathcal{X})$ is the Gaussian width term defined by \eqref{eqn:gauss_width} depending on $\mathcal{X}$. Our fixed design has length $\mathcal{O}\!\left(\frac{d\log(1/\delta)}{\varepsilon^2}+\frac{w(\mathcal{X})^2}{\varepsilon^2}\right)$ and is constructed using two design distributions, $\lambda_1$ and $\lambda_2$. Here, $\lambda_1$ minimizes $\max_{x\in\mathcal{X}}\|x\|^2_{A(\lambda;\mathcal{X})^{-1}}$, while $\lambda_2$ minimizes $\mathbb{E}_{\eta\sim\mathcal{N}(0,I_d)}\!\left[\max_{x\in\mathcal{X}}\big\langle x, A(\lambda;\mathcal{X})^{-1/2}\eta\big\rangle\right]$. Given the resulting samples, we form the unbiased least-squares estimator $\widehat{\theta}$ and output an arm $\widehat{x}\in\arg\max_{x\in\mathcal{X}}\langle x,\widehat{\theta}\rangle$. Using the Borell-TIS inequality, we show that $\widehat{x}$ is an $\varepsilon$-best arm with probability at least $1-\delta$.

We then prove a matching minimax lower bound of $\Omega\!\left(\frac{d\log(1/\delta)}{\varepsilon^2}+\frac{w(\mathcal{X})^2}{\varepsilon^2}\right)$ for any non-adaptive fixed-design algorithm. We first show a lower bound of $\Omega\!\left(\frac{d\log(1/\delta)}{\varepsilon^2}\right)$ for any algorithm (possibly adaptive). For the $\Omega\!\left(\frac{w(\mathcal{X})^2}{\varepsilon^2}\right)$ lower bound, our proof reduces $\varepsilon$-best arm identification to bounding the minimax simple regret after $T$ rounds. We then select a Gaussian prior over the reward vector $\theta$ tailored to the fixed design $(x_t)_{t\in[T]}$, more precisely to the design matrix $\sum_{t=1}^T x_t x_t^\top$, and use it to derive a lower bound on the simple regret. 

We can make our algorithm adaptive as follows. We partition the arm set into $d$ regions and attempt to identify a candidate good arm from each region. We then run Median Elimination on these $d$ candidate arms to obtain an $\varepsilon$-best arm. In certain cases, this yields a $\log d$ improvement, since the analysis can be localized to the region containing the best arm and the Borell--TIS inequality inequality is applied only over that region.

This observation also motivates us to ask whether there are structured action sets beyond the multi-armed bandit case for which adaptivity yields at most logarithmic factors improvement over our non-adaptive algorithm. Towards that we study several structured action sets of interest, namely the hypercube, the unit $\ell_2$-ball, $m$-sets, and multi-task multi-armed bandits (MAB), and characterize the corresponding Gaussian width terms $w(\mathcal{X})$. For each of these sets, we show, by establishing an adaptive lower bound, that adaptive algorithms can improve over our non-adaptive fixed-design sample complexity by at most logarithmic factors only. Our main technical contributions in this part are the lower bound arguments for $m$-sets and multi-task MAB. We begin by constructing a family of hard instances that is oblivious to the algorithm: each instance specifies a reward vector whose coordinates are chosen as $\varepsilon$ times coordinate-specific scaling factor. We then show that, unless the algorithm collects sufficiently many samples, it must incur simple regret $\Omega(\varepsilon)$. To analyze the simple regret, we localize the KL-divergence-based argument. Concretely, we restrict attention to a structured subset of instances in which a portion of the reward vector is fixed and any two instances differ in exactly two coordinates within the remaining portion. This reduction enables a KL-divergence analysis closely analogous to the classical multi-armed bandit setting. Finally, a global averaging argument lifts the localized bound to the full family, yielding simple regret $\Omega(\varepsilon)$ and, consequently, a minimax sample complexity lower bound.

The above result leads us to our main question: does adaptivity in linear bandits improve over non-adaptive approaches by only logarithmic factors? We answer this in the negative by constructing an action set $\mathcal{X}$ for which adaptivity yields a {\it polynomial-factor} improvement over all non-adaptive algorithms. We take $\mathcal{X}$ to be the union of $k=\mathrm{poly}(d)$ unit $\ell_2$-balls, each contained in a distinct $d$-dimensional subspace. A naive non-adaptive approach allocates on the order of $d^2/\varepsilon^2$ samples to each ball to form an unbiased estimator and identify an $\varepsilon$-best arm, leading to total sample complexity $kd^2/\varepsilon^2$. In contrast, an adaptive strategy can save a factor of $d$ by first identifying a ball that contains a near-optimal arm and then allocating additional samples to that ball to recover a near-optimal arm using just on the order of $\frac{kd + d^2}{\varepsilon^2}$ samples.


A key step in identifying such a “good” unit ball is an $\ell_2$-norm estimation subroutine, which aims to estimate the best value over the unit ball, namely the $\ell_2$-norm of the reward vector. Specifically, we design an adaptive algorithm that takes $\mathcal{O}\!\left(\frac{d\log(1/\delta)}{\varepsilon^2}\right)$ samples from the unit $\ell_2$-ball in $\mathbb{R}^d$, observes noisy rewards of the form $\langle x,\theta\rangle+\eta$ for an unknown vector $\theta\in\mathbb{R}^d$, and outputs an estimate $\widehat r$ such that, with probability at least $1-\delta$, $\widehat r\in[\|\theta\|_2-\varepsilon,\|\theta\|_2+\varepsilon]$. A key technical component of our algorithm is to sample multiple Rademacher unit vectors uniformly at random and take multiple observations along each sampled direction. We then construct an $\ell_2$-norm estimator by squaring the empirical mean reward along each direction and aggregating these squared estimates. This squaring step leads to a subexponential-tail analysis, which is the main non-trivial part of our proof required to establish the stated sample complexity. 

By constructing such a set $\mathcal{X}$, we highlight a key technical insight: accurately estimating the best value over a region $\mathcal{Z}\subset\mathcal{X}$, namely $\max_{x\in\mathcal{Z}}\langle x,\theta\rangle$, plays a crucial role in achieving improved performance. In the special case where $\mathcal{Z}$ is a unit $\ell_2$-ball, this task reduces to $\ell_2$-norm estimation. More generally, this observation suggests that region-wise value estimation is a fundamental ingredient in the design of near-optimal policies in related settings.

% \textcolor{red}{Maybe emphasize that the approach for this set presents a departure from the standard approach. Also stress this is a key technical contribution suggesting that future works should take a different approach. Maybe good for conclusion section.}
\subsection{Related Works}
Our work is closely related to a line of research on instance-dependent sample complexity for best arm identification \citep{soare2014best,karnin2016verification,xu2018fully,tao2018best,fiez2019sequential,jedra2020optimal,katz2020empirical,degenne2020gamification}. A key distinction between instance-dependent and minimax guarantees is the role of $\theta$. In the instance-dependent regime, one seeks sample complexity bounds of the form $\mathbf{H}_1\log(1/\delta)+\mathbf{H}_2$, where $\mathbf{H}_1$ and $\mathbf{H}_2$ depend on the arm set $\mathcal{X}$ and the reward vector $\theta$. In contrast, minimax characterizations aim for bounds of the same form with $\mathbf{H}_1$ and $\mathbf{H}_2$ depending on the arm set $\mathcal{X}$ and the accuracy parameter $\varepsilon$, but not on $\theta$, since the guarantee is worst-case over $\theta$. While several instance-dependent algorithmic ideas extend to the minimax setting, lower bounds are less transferable since they typically hinge on simple-regret arguments, and minimax complexity characterizations remain relatively sparse \citep{even2002pac,shamir2015complexity,chen2024assouad}.


One approach to designing algorithms in the instance-dependent regime is through experimental design. Early works such as \cite{soare2014best} and \cite{karnin2016verification} used G-optimal designs. Later, \cite{fiez2019sequential} proposed an algorithm based on sequential experimental design that is near-optimal for the first term $\mathbf{H}_1$, but can incur a large second term $\mathbf{H}_2$. This motivated \cite{katz2020empirical} to design an algorithm whose design distribution minimizes certain Gaussian width terms, thereby reducing $\mathbf{H}_2$. Such a Gaussian-width based approach to experimental design is also used in a related regret minimization setting by \cite{wagenmaker2021experimental}. For a detailed discussion of experimental design and its connection to linear bandits, we refer the reader to \cite{lattimore2020bandit}.

Our work is also related to the literature on norm estimation and value estimation. \cite{kong2020sublinear} studied optimal policy value estimation in contextual bandits. In a different direction, \cite{cai2011testing} studied estimating $\frac{1}{n}\sum_i |\theta_i|$ from an observation $Y\sim\mathcal{N}(\theta,I_n)$, and \cite{collier2020estimation} later generalized this to estimating $\frac{1}{n}\sum_i |\theta_i|^\gamma$ for $\gamma>0$. \cite{han2020estimation} studied $L_r$-norm estimation in Gaussian white noise models. More recently, \cite{cleanthous2025adaptive} studied adaptive estimation of the $L_2$-norm of a probability density on $\mathbb{R}^d$.

